\documentclass[11pt]{article}

\usepackage{amsmath,amssymb,amsthm}
\usepackage{hyperref}
\usepackage{geometry}
\geometry{margin=1in}

\title{Half-energy Page time from dimension balance in energy-constrained quantum systems}
\author{}
\date{}

\newtheorem{theorem}{Theorem}
\newtheorem{proposition}{Proposition}
\newtheorem{remark}{Remark}

\begin{document}
\maketitle

\begin{abstract}
Page-curve--like non-monotonic entropy dynamics has been observed in a variety of
quantum systems, including strictly Markovian cooling models, where the entropy
turnover coincides with emission of approximately half the initial energy.
We provide a general explanation for this coincidence using a constrained-random-state
(Page/typicality) framework.
Within this framework, entropy turnover occurs when the effective Hilbert-space
dimensions of the accessible energy sectors of a system and its complement become
comparable.
We show that when the relevant microcanonical entropies are approximately linear in
energy---equivalently, when the density of states is exponential with a common slope---
this dimension-balance condition generically yields a half-energy turnover, up to
subextensive corrections.
For more general entropy--energy relations, the turnover fraction is model-dependent.
\end{abstract}

\section{Introduction}

The Page curve---the characteristic rise and fall of subsystem entropy in a globally
pure quantum state---was originally derived for random pure states on bipartite
Hilbert spaces and depends primarily on relative subsystem dimensions \cite{Page1993}.
Page-curve--like entropy dynamics has since been observed in a wide range of settings,
including interacting many-body systems, open quantum systems, and thermodynamic
relaxation models.

In particular, Glatthard studied strictly Markovian Lindblad dynamics describing energy
relaxation to a cold environment and observed Page-curve--like entropy turnover, with
the notable feature that the turnover occurs when approximately half of the initial
energy has been emitted \cite{Glatthard2025}. Similar behavior has been reported in
analytically solvable free-fermion models \cite{Kehrein2024} and in interacting XXZ
spin chains coupled to baths \cite{RayDharKulkarni2025}.

The purpose of this note is to explain when and why such half-energy turnover arises.
Our analysis does not rely on non-Markovianity, gravity, or black-hole-specific
dynamics. Instead, it rests on a balance condition between energy-resolved Hilbert-space
dimensions, closely analogous to Page’s original argument.

\section{Framework and dimension-balance condition}

Consider a bipartite quantum system $A+B$ with Hamiltonian $H=H_A+H_B$.
Fix a total energy $E_0$ and an energy window defining a microcanonical shell.
Let $\Omega_X(E)$ denote the number of states of subsystem $X\in\{A,B\}$ in an energy
window around $E$, and define the corresponding microcanonical entropies
\[
S_X(E) := \log \Omega_X(E).
\]

\paragraph{Framework assumption.}
We assume that pure states drawn uniformly from the global energy shell exhibit
canonical typicality: with high probability, reduced states are close to
microcanonical or canonical states determined by the energy constraint
\cite{PopescuShortWinter2006,Goldstein2006}.
Within this framework, entanglement is controlled by the smaller effective
Hilbert-space dimension compatible with the constraint.

\paragraph{Dimension-balance condition.}
Under the assumption above, entropy turnover occurs when the effective dimensions of
the two subsystems become comparable, leading to the balance condition
\begin{equation}
\label{eq:balance}
S_A(E_A^\star) = S_B(E_0 - E_A^\star),
\end{equation}
up to finite-size and energy-window corrections.
Equation \eqref{eq:balance} is the energy-resolved analogue of Page’s
dimension-balance criterion for random pure states \cite{Page1993}.

\paragraph{Illustrative example (no dynamics).}
If $\Omega_A(E)\propto e^{\beta E}$ and $\Omega_B(E)\propto e^{\beta E}$ over the
relevant range, then $S_A(E)=\beta E+\mathrm{const}$ and
$S_B(E)=\beta E+\mathrm{const}$, and \eqref{eq:balance} yields
$E_A^\star=E_0/2$.
This example is intended solely to visualize the balance condition and does not rely
on any specific dynamical model.

\section{Half-energy turnover from exponential density of states}

\begin{theorem}[Half-energy turnover]
\label{thm:half_energy}
Assume that over the relevant energy range the microcanonical entropies satisfy
\[
S_A(E)=\beta E + c_A + r_A(E), \qquad
S_B(E)=\beta E + c_B + r_B(E),
\]
with a common slope $\beta>0$ and subextensive corrections
$r_A(E),r_B(E)=o(E_0)$ uniformly for $E\in[0,E_0]$.
Then any solution of \eqref{eq:balance} satisfies
\[
E_A^\star = \frac{E_0}{2} + \frac{c_B-c_A}{2\beta} + o(E_0),
\]
and in the thermodynamic limit $E_A^\star=E_0/2+o(E_0)$.
\end{theorem}

\begin{proof}
Substituting the assumed forms into \eqref{eq:balance} gives
\[
\beta E_A^\star + c_A + r_A(E_A^\star)
=
\beta(E_0-E_A^\star) + c_B + r_B(E_0-E_A^\star).
\]
Solving for $E_A^\star$ yields the stated result at leading extensive order, provided
the subextensive terms do not generate an $O(E_0)$ shift.
\end{proof}

\section{Non-half-energy turnover for general entropies}

\begin{proposition}[Power-law entropies]
\label{prop:powerlaw}
Assume the dimension-balance condition \eqref{eq:balance} holds and that
$S_A(E)=\alpha_A E^\gamma$ and $S_B(E)=\alpha_B E^\gamma$ with
$\alpha_{A,B}>0$ and $\gamma>0$.
Then
\[
E_A^\star = \frac{E_0}{1+(\alpha_A/\alpha_B)^{1/\gamma}}.
\]
If $\alpha_A=\alpha_B$, then $E_A^\star=E_0/2$ for any $\gamma$.
This symmetry-based half-energy result is distinct from the linear-entropy mechanism
of Theorem~\ref{thm:half_energy}.
\end{proposition}

\section{Relation to Markovian cooling and many-body models}

The half-energy turnover observed in strictly Markovian cooling models
\cite{Glatthard2025} fits naturally into the framework above: the turnover reflects a
balance of accessible Hilbert-space dimensions during steady energy transfer, rather
than information backflow.
Related Page-curve--like behavior has been reported in free-fermion models
\cite{Kehrein2024} and in interacting spin chains coupled to baths
\cite{RayDharKulkarni2025}.

\section{Remark on black holes and scope}

\begin{remark}[Black holes]
For a Schwarzschild black hole, the Bekenstein--Hawking entropy scales as
$S_{\rm BH}(M)\propto M^2$, so the density of states behaves as
$\Omega(M)\sim e^{M^2}$.
If the combined black hole--radiation system evolves unitarily and remains in a pure
state, then $S_{\rm rad}=S_{\rm BH}^{(\mathrm{red})}$ at all times.
Under the additional assumption that the total entropy is conserved and initially
resides entirely in the black hole, one has
$S_{\rm rad}\approx S_{\rm BH}(M_0)-S_{\rm BH}(M)$.
The Page-time condition $S_{\rm rad}\approx S_{\rm BH}(M)$ then yields
$S_{\rm BH}(M)\approx \tfrac12 S_{\rm BH}(M_0)$, corresponding to
$M_{\rm Page}\approx M_0/\sqrt{2}$.
Black holes therefore realize a different entropy--energy relation together with
explicit unitarity assumptions, leading to a different turnover fraction.
\end{remark}

\section{Conclusion}

Half-energy Page times arise generically from a simple dimension-balance condition
when the relevant microcanonical entropies are approximately linear in energy.
When this linearity fails, the turnover fraction is model-dependent.
These results clarify the structural origin of half-energy entropy turnover observed
in Markovian cooling models and delineate the limits of its applicability.

\bibliographystyle{unsrt}
\bibliography{half_energy_page_time}

\end{document}
